% ============================================================
\section{Related Work}\label{sec:related}
% ============================================================

Six research streams converge on the problem the Cascade addresses.

\paragraph{The structural hole at the epistemic boundary.} Recent technical results specify the governance gap from complementary angles. Nikolaou et al.~\cite{nikolaou2025injective} establish that, under their assumptions, language-model mappings can be injective and hence invertible. Activation-engineering work shows that model behavior can also be steered by direct activation interventions rather than by textual prompting alone~\cite{turner2024steering}. Mechanistic-interpretability surveys catalog the progress and limits of circuit- and representation-level explanations for alignment~\cite{naseem2026mechanistic}. \textit{Dual Certificates}~\cite{anon2026dual} is the companion manuscript from which this paper imports its stronger information-theoretic premise. Together, these sources motivate the epistemic-boundary claim; the proof obligation for the full non-recoverability statement remains in the companion result.

\paragraph{Layered governance architectures.} Several recent proposals organize AI governance into layered structures~\cite{bommarito2025governing}. These are \textit{regulatory} architectures specifying what rules should apply at each level; the Cascade is a \textit{diagnostic} architecture identifying where damage originates and how it propagates---the question regulatory approaches presuppose is answered. A regulatory-layer analysis of the Garcia case would say ``product liability applies'' but cannot identify \textit{which architectural component} produced the harm or predict cross-layer propagation.

\paragraph{Product liability and the classification gap.} Ayres \& Balkin~\cite{ayres2024risky} establish that AI systems lack intentions and should be held to objective liability standards. Abraham \& Sharkey~\cite{abraham2026untangling} provide a comprehensive tort analysis, identifying the ``black box'' problem as a core barrier to doctrinal fit. Ramakrishnan et al.~\cite{ramakrishnan2024tortai} survey how tort exposure for large-scale AI harms varies across developers, deployers, and users. Selbst et al.'s ``Ripple Effect Trap''~\cite{selbst2019fairness}---where interventions at one system level produce unforeseen consequences at others---is the motivating insight the Cascade formalizes. In the medical setting specifically, the proposed EU liability directives are found to leave gaps that fall on precisely this black-box seam~\cite{duffourc2023}, and common-enterprise theories have been advanced to reallocate the liability those gaps orphan~\cite{chan2021}---doctrinal moves that presuppose the architectural diagnosis the Cascade supplies.

\paragraph{Mens rea impossibility.} Abbott \& Sarch~\cite{abbott2020punishing} establish through doctrinal analysis that \textit{mens rea} attribution is structurally impossible for AI systems. Nerantzi \& Sartor's ``hard AI crime''~\cite{nerantzi2024hardai} identifies the resulting criminal responsibility gap across jurisdictions. Alternative approaches propose new legal fictions: Mukherjee \& Chang's ``operational agency''~\cite{mukherjee2026operational} would treat AI as possessing permeable agency for culpability-tracing. The Cascade avoids this jurisprudential cost: rather than extending legal personhood concepts, it diagnoses which architectural layer failed.

\paragraph{Single-layer documentation and the verification gap.} Model Cards~\cite{mitchell2019model} and Datasheets~\cite{gebru2021datasheets} established transparency standards; HELM~\cite{liang2023holistic} benchmarks cross-scenario robustness. Benerofe~\cite{benerofe2025verification} argues that computational intractability creates a structural verification gap. Each operates within a single layer; individually verified layers can still produce emergent damage through interaction. Imana et al.~\cite{imana2025meta}---FAccT 2025 Best Paper---demonstrated this concretely: Meta's Variance Reduction System, an L4 intervention, achieved demographic parity in ad delivery by suppressing reach for all groups rather than expanding it for underserved groups---an output-policing fix that reduced total utility without improving individual opportunity.

\paragraph{Systems safety engineering.} Leveson's \textit{Engineering a Safer World}~\cite{leveson2012engineering} and the STPA Handbook~\cite{leveson2018stpa} established the gold standard for layered accident analysis. Rismani et al.~\cite{rismani2024silos} propose process-oriented hazard analysis to overcome the ``silo problem'' in AI safety. The Cascade inherits this layered diagnostic architecture but operates where STPA's prerequisite fails: AI deployment lacks the assumption that operational safety decomposes into well-defined physical sub-states. Terminal outcomes may be uncontested while operational safe states remain disputed---the gap the Cascade addresses. Madaio et al.~\cite{madaio2024rai}---FAccT 2024 Best Paper---provide qualitative evidence: practitioners report learning responsible AI through on-the-job friction rather than formal training, confirming that safety knowledge itself propagates through the same lossy cascade this paper diagnoses.

